\section{Nền tảng lý thuyết}
\label{sec:theoretical-background}

\subsection{Bài toán khai thác tập mục phổ biến}

Khai thác tập mục phổ biến (\textit{Frequent Itemset Mining}) là bài toán
tìm các tập mục xuất hiện với tần suất không nhỏ hơn một ngưỡng cho trước trong
cơ sở dữ liệu giao dịch. Đây là nền tảng của khai phá luật kết hợp và được sử
dụng trong nhiều bài toán như phân tích giỏ hàng, phân tích hành vi, khai phá log
và phát hiện các mẫu đồng xuất hiện~\cite{agrawal1994fast,han2011datamining}.

Cho tập tất cả các mục:
\begin{equation}
    I = \{i_1,i_2,\ldots,i_m\}.
    \label{eq:item-universe}
\end{equation}

Một \textit{itemset} là một tập con của $I$. Cơ sở dữ liệu giao dịch
(\textit{transaction database}, TDB) được ký hiệu:
\begin{equation}
    D = \{T_1,T_2,\ldots,T_n\},
    \label{eq:transaction-database}
\end{equation}
trong đó mỗi giao dịch $T_j \subseteq I$ là tập các mục xuất hiện cùng nhau
trong một lần quan sát và thường được gắn với một mã giao dịch
(\textit{transaction identifier}, TID).

Với itemset $X \subseteq I$, giao dịch $T_j$ chứa $X$ khi và chỉ khi:
\begin{equation}
    X \subseteq T_j.
    \label{eq:transaction-contains-itemset}
\end{equation}

Độ hỗ trợ tuyệt đối của $X$ trên $D$ là số giao dịch chứa $X$:
\begin{equation}
    \operatorname{sup}_{abs}(X)
    =
    \left|\left\{T_j \in D \mid X \subseteq T_j\right\}\right|.
    \label{eq:absolute-support}
\end{equation}

Độ hỗ trợ tương đối là tỉ lệ giao dịch chứa $X$:
\begin{equation}
    \operatorname{sup}_{rel}(X)
    =
    \frac{\operatorname{sup}_{abs}(X)}{|D|}.
    \label{eq:relative-support}
\end{equation}

Trong báo cáo này, nếu không nêu rõ, $\operatorname{sup}(X)$ được hiểu là độ
hỗ trợ tuyệt đối. Với ngưỡng hỗ trợ tối thiểu $\mathit{minsup}$ được biểu diễn
theo số giao dịch, $X$ là một tập mục phổ biến (\textit{frequent itemset})
khi:
\begin{equation}
    \operatorname{sup}(X) \geq \mathit{minsup}.
    \label{eq:frequent-itemset-condition}
\end{equation}

Do đó, bài toán FIM được phát biểu:
\begin{equation}
    FI(D,\mathit{minsup})
    =
    \left\{
        X \subseteq I
        \mid
        \operatorname{sup}(X) \geq \mathit{minsup}
    \right\}.
    \label{eq:fim-problem}
\end{equation}

Bên cạnh FI thông thường, hai dạng biểu diễn rút gọn thường được xét là tập mục
phổ biến đóng (\textit{closed frequent itemset}) và tập mục phổ biến cực
đại (\textit{maximal frequent itemset})~\cite{han2011datamining}.

Một FI $X$ là CFI nếu không tồn tại FI $Y$ thỏa:
\begin{equation}
    X \subset Y
    \quad\text{và}\quad
    \operatorname{sup}(X)=\operatorname{sup}(Y).
    \label{eq:closed-frequent-itemset}
\end{equation}

Một FI $X$ là MFI nếu không tồn tại FI $Y$ thỏa:
\begin{equation}
    X \subset Y.
    \label{eq:maximal-frequent-itemset}
\end{equation}

Quan hệ giữa ba tập kết quả là:
\begin{equation}
    MFI \subseteq CFI \subseteq FI.
    \label{eq:fi-cfi-mfi-relation}
\end{equation}

MFI cho kết quả gọn hơn nhưng không bảo toàn trực tiếp support của mọi tập con.
CFI vẫn là biểu diễn không dư thừa theo quan hệ support và có thể được dùng để
suy ra support của các FI thông thường.

\subsection{Tính chất Apriori}

Tính chất Apriori, còn gọi là tính chất phản đơn điệu
(\textit{anti-monotonicity}) của support, phát biểu rằng với hai itemset
$X \subseteq Y$:
\begin{equation}
    \operatorname{sup}(Y) \leq \operatorname{sup}(X).
    \label{eq:apriori-property}
\end{equation}

Thật vậy, mọi giao dịch chứa $Y$ đều chứa $X$, nên:
\begin{equation}
    \left\{T_j \in D \mid Y \subseteq T_j\right\}
    \subseteq
    \left\{T_j \in D \mid X \subseteq T_j\right\}.
    \label{eq:apriori-set-inclusion}
\end{equation}

Từ đó thu được hai hệ quả quan trọng:
\begin{itemize}
    \item Nếu $X$ không phổ biến thì mọi tập cha $Y \supset X$ cũng không phổ biến.
    \item Nếu $Y$ phổ biến thì mọi tập con $X \subseteq Y$ đều phổ biến.
\end{itemize}

Apriori sử dụng trực tiếp tính chất này để sinh và loại ứng viên theo từng mức
kích thước~\cite{agrawal1994fast}. Các thuật toán về sau như Eclat, FP-Growth
và PrePost vẫn dựa trên tính phản đơn điệu để tỉa không gian tìm kiếm, nhưng sử
dụng các cấu trúc biểu diễn và chiến lược khai thác khác nhau
~\cite{zaki2000eclat,han2000fpgrowth,deng2012prepost}.

\subsection{Tổng quan thuật toán PrePost}

PrePost là thuật toán khai thác toàn bộ FI dựa trên cấu trúc
\textit{Pre-order Post-order Coding tree} và danh sách mã hóa node ~\cite{deng2012prepost}. Thuật toán được thiết kế để tránh việc
quét lại TDB trong giai đoạn mở rộng itemset và giảm chi phí biểu diễn support
so với việc lưu trực tiếp danh sách TID dài.

Quy trình tổng quát gồm:
\begin{enumerate}
    \item Quét TDB để tính support của từng item và loại các item không phổ biến.
    \item Sắp xếp các item phổ biến theo một thứ tự toàn cục, thường là support
          giảm dần.
    \item Chèn các giao dịch đã lọc và sắp xếp vào PPC-tree.
    \item Duyệt cây để gán mã \textit{pre-order} và \textit{post-order}.
    \item Xây dựng N-list cho các frequent 1-itemset.
    \item Mở rộng itemset theo chiều sâu bằng phép giao N-list và loại các
          itemset có support nhỏ hơn $\mathit{minsup}$.
\end{enumerate}

Với hai node $u$ và $v$ trên PPC-tree, $u$ là tổ tiên của $v$ khi:
\begin{equation}
    pre(u) < pre(v)
    \quad\text{và}\quad
    post(u) > post(v).
    \label{eq:ancestor-condition}
\end{equation}

Điều kiện này cho phép kiểm tra quan hệ tổ tiên--hậu duệ chỉ bằng hai số nguyên,
không cần duyệt lại đường đi trên cây. Đây là cơ sở để PrePost tính support của
các itemset mở rộng thông qua N-list~\cite{deng2012prepost}.

\subsection{Cấu trúc PPC-tree}

PPC-tree là cây tiền tố được xây dựng từ TDB sau khi loại các item không phổ
biến và sắp xếp các item còn lại theo thứ tự toàn cục. Việc đặt các item có
support cao ở gần gốc thường làm tăng khả năng chia sẻ tiền tố giữa các giao
dịch, từ đó giảm số node cần lưu.

Mỗi node PPC-tree lưu:
\begin{itemize}
    \item \textbf{item-name}: item được biểu diễn bởi node;
    \item \textbf{count}: số giao dịch đi qua tiền tố kết thúc tại node;
    \item \textbf{children}: tập hoặc danh sách node con;
    \item \textbf{pre}: thứ tự node khi được thăm trong phép duyệt trước;
    \item \textbf{post}: thứ tự node khi hoàn tất trong phép duyệt sau.
\end{itemize}

Mỗi node được gán một cặp mã:
\begin{equation}
    PPCode(u)=(pre(u),post(u)).
    \label{eq:pp-code}
\end{equation}

Các giá trị $pre$ và $post$ phải được sinh nhất quán bởi cùng một quá trình
duyệt cây. Giá trị cụ thể không quan trọng bằng việc chúng bảo toàn điều kiện
tổ tiên trong Công thức~\eqref{eq:ancestor-condition}.

\subsection{Cấu trúc N-list}

N-list của một item hoặc itemset là một danh sách các bộ ba:
\begin{equation}
    (pre,post,count).
    \label{eq:nlist-entry}
\end{equation}

Với frequent 1-itemset $\{i\}$, $NList(i)$ chứa PP-code và count của tất cả node
mang nhãn $i$. Support được tính bằng:
\begin{equation}
    \operatorname{sup}(i)
    =
    \sum_{e \in NList(i)} e.count.
    \label{eq:nlist-single-support}
\end{equation}

Với itemset dài hơn, N-list được xây dựng từ N-list của hai tiền tố có quan hệ
phù hợp trong cây liệt kê tập hợp. Phép kết hợp chỉ giữ các cặp node thỏa quan
hệ tổ tiên--hậu duệ theo mã pre/post. Khi nhiều phần tử tạo ra cùng một PP-code
đầu ra, count của chúng phải được cộng gộp. Có thể ký hiệu phép này dưới dạng:
\begin{equation}
    NList(P \cup \{x\})
    =
    NList(P) \bowtie_N NList(x),
    \label{eq:nlist-join}
\end{equation}
trong đó $\bowtie_N$ là phép giao N-list theo định nghĩa và thứ tự item của
PrePost~\cite{deng2012prepost}.

Sau khi giao, support của itemset được tính trực tiếp từ N-list:
\begin{equation}
    \operatorname{sup}(P \cup \{x\})
    =
    \sum_{e \in NList(P \cup \{x\})} e.count.
    \label{eq:nlist-itemset-support}
\end{equation}

Điểm cần lưu ý là N-list của một itemset không nhất thiết chỉ có một phần tử.
Số phần tử phụ thuộc vào số node PPC-tree thỏa quan hệ cần thiết; các phần tử
chỉ được hợp nhất khi chúng có cùng PP-code đầu ra.

Nhờ N-list, sau khi PPC-tree được xây dựng, thuật toán không cần quét lại TDB để
tính support cho từng itemset mở rộng. Phần lớn chi phí được chuyển thành các
phép duyệt và giao trên các danh sách mã hóa node~\cite{deng2012prepost}.

\subsection{Giả mã thuật toán PrePost}

Thuật toán~\ref{alg:prepost} mô tả quy trình PrePost ở mức tổng quát. Giả mã
tách rõ hai giai đoạn: tiền xử lý để xây PPC-tree và N-list, sau đó khai thác
theo chiều sâu.

\begin{lstlisting}[
    caption={Giả mã tổng quát của thuật toán PrePost},
    label={alg:prepost}
]
Input : Transaction database D, minimum support minsup
Output: All frequent itemsets FI

1. Scan D and compute the support of every item.
2. Remove items having support < minsup.
3. Build a global order F of frequent items by descending support.
4. Create an empty PPC-tree.
5. For each transaction T in D:
       a. Remove infrequent items from T.
       b. Sort the remaining items according to F.
       c. Insert the ordered transaction into the PPC-tree.
6. Assign pre-order and post-order codes to all PPC-tree nodes.
7. Build the N-list of every frequent 1-itemset.
8. Call Mine(prefix = empty, extensions = frequent 1-itemsets).

Procedure Mine(prefix, extensions):
1. For each extension X in extensions:
       a. P = prefix union {X}.
       b. Output P.
       c. next_extensions = empty.
       d. For each extension Y occurring after X:
              i.   Construct NList(P union {Y}) by N-list intersection.
              ii.  Compute support from the resulting N-list.
              iii. If support >= minsup:
                       store the N-list;
                       add Y to next_extensions.
       e. If next_extensions is not empty:
              Mine(P, next_extensions).
\end{lstlisting}

Thủ tục \texttt{Mine} duyệt không gian tìm kiếm theo chiều sâu. Một itemset mở
rộng chỉ được đưa vào mức đệ quy tiếp theo khi support của nó không nhỏ hơn
$\mathit{minsup}$. Việc loại sớm này là hệ quả trực tiếp của tính phản đơn điệu:
nếu itemset hiện tại không phổ biến thì mọi phần mở rộng của nó đều không phổ
biến.

\subsection{Phân tích độ phức tạp}

Hiệu năng của PrePost phụ thuộc vào số giao dịch, số item phân biệt, độ dài giao
dịch, ngưỡng $\mathit{minsup}$, mức độ chia sẻ tiền tố và số FI trong kết quả.

Gọi $n=|D|$, $m=|I|$ và $\ell$ là độ dài trung bình của giao dịch sau lọc.
Quét TDB để đếm support của item đơn có chi phí:
\begin{equation}
    O(n\ell).
    \label{eq:scan-complexity}
\end{equation}

Nếu các item trong mỗi giao dịch cần được sắp xếp độc lập, chi phí tiền xử lý
có thể bao gồm:
\begin{equation}
    O\!\left(
        \sum_{T \in D}|T|\log |T|
    \right).
    \label{eq:transaction-sort-complexity}
\end{equation}

Sau khi giao dịch đã được sắp xếp, việc chèn vào PPC-tree xử lý mỗi lần xuất
hiện item ít nhất một lần. Khi phép tra cứu node con được cài đặt bằng bảng băm
hoặc ánh xạ phù hợp, chi phí xây cây thường gần:
\begin{equation}
    O(n\ell).
    \label{eq:tree-building-complexity}
\end{equation}

Tuy nhiên, không thể biểu diễn toàn bộ độ phức tạp khai thác chỉ bằng $n$, $m$
và $\ell$, vì chi phí còn phụ thuộc vào số itemset được xét và độ dài các
N-list. Nếu một phép giao hai N-list được cài đặt bằng hai con trỏ trên các
danh sách đã sắp xếp, chi phí của phép giao thường tỉ lệ với tổng độ dài hai
danh sách:
\begin{equation}
    O\!\left(
        |NList(X)|+|NList(Y)|
    \right).
    \label{eq:nlist-intersection-complexity}
\end{equation}

Tổng thời gian khai thác có thể mô tả ở mức tổng quát:
\begin{equation}
    O\!\left(
        \sum_{c \in \mathcal{C}}
        \operatorname{costJoin}(c)
    \right),
    \label{eq:mining-total-complexity}
\end{equation}
trong đó $\mathcal{C}$ là tập các phép mở rộng thực sự được thuật toán xét.

Trong trường hợp xấu nhất, số FI có thể tăng theo hàm mũ theo số item. Đây là
giới hạn đầu ra của bài toán FIM chứ không chỉ riêng PrePost. Vì vậy, khi dữ
liệu dày đặc và $\mathit{minsup}$ thấp, thời gian chạy và bộ nhớ có thể tăng
mạnh do số itemset và N-list trung gian rất lớn
~\cite{han2011datamining,luna2019review}.

Độ phức tạp không gian gồm:
\begin{itemize}
    \item bộ nhớ lưu PPC-tree;
    \item bộ nhớ lưu N-list của frequent 1-itemset;
    \item bộ nhớ cho các N-list trung gian trên các nhánh khai thác đang hoạt động;
    \item bộ nhớ lưu kết quả, nếu toàn bộ FI được giữ trong RAM.
\end{itemize}

Các yếu tố ảnh hưởng chính đến hiệu năng là:
\begin{enumerate}
    \item \textbf{Ngưỡng $\mathit{minsup}$}: ngưỡng thấp thường tạo nhiều FI và
          nhiều phép giao N-list hơn.
    \item \textbf{Độ dày dữ liệu}: dữ liệu dày đặc có xu hướng tạo nhiều
          itemset dài.
    \item \textbf{Độ dài giao dịch}: giao dịch dài làm tăng chi phí xây cây và
          số khả năng kết hợp.
    \item \textbf{Mức độ chia sẻ tiền tố}: chia sẻ cao giúp PPC-tree nén tốt hơn.
    \item \textbf{Cài đặt N-list}: cách tổ chức dữ liệu, giao danh sách, cộng gộp
          count và quản lý cấp phát ảnh hưởng trực tiếp đến thời gian chạy.
    \item \textbf{Chiến lược lưu ứng viên}: loại ngay itemset không đạt
          $\mathit{minsup}$ giúp tránh lưu và mở rộng các nhánh chắc chắn thất bại.
\end{enumerate}

\subsection{So sánh với các thuật toán FIM tiêu biểu}

\paragraph{Apriori.}
Apriori khai thác theo chiều rộng và sinh candidate $k$-itemset từ các
frequent $(k-1)$-itemset. Thuật toán dễ hiểu và dễ kiểm chứng, nhưng có thể tạo
rất nhiều ứng viên và cần quét TDB qua nhiều mức
~\cite{agrawal1994fast,han2011datamining}.

\paragraph{Eclat.}
Eclat sử dụng biểu diễn dọc, trong đó mỗi item hoặc itemset được gắn với một
tidset. Support của itemset mở rộng được tính bằng phép giao tidset. Phương pháp
này phù hợp với khai thác theo chiều sâu và tránh quét lại TDB, nhưng các tidset
có thể lớn khi số giao dịch cao~\cite{zaki2000eclat}.

\paragraph{FP-Growth.}
FP-Growth xây dựng FP-tree để nén TDB và khai thác mẫu bằng chiến lược tăng
trưởng mẫu, tránh sinh candidate theo từng mức như Apriori
~\cite{han2000fpgrowth}. Trong quá trình khai thác, thuật toán sử dụng cơ sở mẫu
điều kiện và các FP-tree điều kiện; chi phí của các cấu trúc điều kiện phụ thuộc
vào đặc trưng dữ liệu.

\paragraph{PrePost.}
PrePost cũng sử dụng cây tiền tố, nhưng sau khi xây PPC-tree, thuật toán mã hóa
các node bằng pre/post và biểu diễn itemset bằng N-list. Việc mở rộng itemset
được thực hiện qua phép giao N-list thay vì xây dựng lặp lại các cây điều kiện
~\cite{deng2012prepost}. So với Apriori, PrePost tránh sinh ứng viên theo từng
mức và tránh quét TDB nhiều lần. So với Eclat, PrePost thay tidset bằng biểu
diễn node nén theo PPC-tree. So với FP-Growth, PrePost chuyển phần lớn giai đoạn
khai thác sang thao tác danh sách.

PrePost không luôn vượt trội trên mọi loại dữ liệu. Khi dữ liệu ít chia sẻ tiền
tố hoặc $\mathit{minsup}$ rất thấp, PPC-tree và các N-list trung gian vẫn có thể
lớn. Hiệu năng thực tế còn phụ thuộc mạnh vào cấu trúc dữ liệu và cách cài đặt
phép giao N-list.

PrePost+ tiếp tục cải tiến PrePost bằng chiến lược
\textit{Children--Parent Equivalence pruning} để giảm số nhánh cần khai thác
~\cite{deng2015prepostplus}. Trong phần thực nghiệm của báo cáo, SPMF được dùng
như một triển khai tham chiếu vì đây là thư viện Java mã nguồn mở chuyên về
khai phá mẫu và cung cấp nhiều thuật toán FIM
~\cite{fournier2014spmf,spmfwebsite}. Các bộ dữ liệu chuẩn được lấy từ FIMI
Repository nhằm hỗ trợ so sánh trên những tập dữ liệu thường dùng trong nghiên
cứu FIM~\cite{fimirepository,goethals2004fimi}.

\subsection{Kết luận chương}

Chương này đã trình bày bài toán FIM, các khái niệm support, FI, CFI, MFI và
tính chất Apriori. Trên cơ sở đó, chương giới thiệu thuật toán PrePost với hai
cấu trúc cốt lõi là PPC-tree và N-list, mô tả quy trình khai thác, giả mã và các
yếu tố chi phối độ phức tạp. Phần so sánh cho thấy PrePost thuộc nhóm thuật toán
khai thác theo chiều sâu, tránh quét lại TDB bằng cách mã hóa thông tin support
trên cây và thực hiện các phép giao N-list. Đây là nền tảng để chương tiếp theo
minh họa quá trình xây PPC-tree, gán mã pre/post và sinh N-list trên một ví dụ
cụ thể.
